% GENERATED FILE. Edit canonical lessons, then run npm run generate:books.
% canonical-source-hash: fnv1a64:6b4d51f3d6d212e6
% canonical-lessons: PA-C207-hona-return, PA-C207-hundaa, PA-C207-hundaa-review, PA-C207-daily-frame, PA-C207-payoff

\chapter{This Happens Every Day}
\label{ch:pa-a2-hunda}
\hlchaptermodality{207}

\begin{chapteropening}
\textbf{By the end of this chapter:} \emph{I can recognise and use one bounded masculine singular habitual form for happens.}

Recognise and produce one bounded this-happens-every-day frame through listening, speaking, reading, and writing.
\end{chapteropening}

\section[\texorpdfstring{\pa{ਹੋਣਾ} · \pa{ਰੋਜ਼} (hona · roz)}{ਹੋਣਾ · ਰੋਜ਼ (hona · roz)}]{Return to to be and every day}

\label{lesson:PA-C207-hona-return}



\begin{quote}
\emph{Hear:} \emph{hona, roz}; YOU SAY: \textbf{\pa{ਹੋਣਾ} · \pa{ਰੋਜ਼}}
\end{quote}

\begin{grammarlens}[title={Grammar Lens}]
Keep the familiar infinitive “to be” and the familiar time word “every day” ready. The next lesson will own one new form; this opener does not reveal it.
\end{grammarlens}

\subsection*{Your turn}
\begin{itemize}
\raggedright
  \item \emph{Hear:} \emph{to be, every day}; YOU SAY: \textbf{\pa{ਹੋਣਾ} · \pa{ਰੋਜ਼}}
  \item Say the two familiar items once, then stop.
\end{itemize}

\begin{tcolorbox}[breakable,colback=teal!4,colframe=teal!35!black,title={Before you move on}]
Produce the familiar infinitive and time word from meaning. This opener owns no new form.
\end{tcolorbox}
\section[\texorpdfstring{\pa{ਹੁੰਦਾ} (hundaa)}{ਹੁੰਦਾ (hundaa)}]{One bounded habitual form of to be}

\label{lesson:PA-C207-hundaa}



\begin{quote}
Retrieve familiar \textbf{\pa{ਹੋਣਾ}}, “to be,” then listen for one masculine singular habitual form.
\end{quote}

\begin{grammarlens}[title={Grammar Lens: direct owner}]
\begin{quote}
\textbf{\pa{ਹੁੰਦਾ}} — \emph{hundaa} — happens; masculine singular habitual form
\end{quote}

Hear and say the whole form before reading it. Keep \textbf{\pa{ਹੁੰਦਾ}} whole. This chapter uses it only for one repeated event in the small frame that follows. It does not open the wider habitual agreement system.
\end{grammarlens}

\begin{cousinweb}
Punjabi University Patiala's lexical statistics identifies \textbf{\pa{ਹੁੰਦਾ}} as the masculine singular non-perfect form of \textbf{\pa{ਹੋ}}, with continuous or habitual use. Its reference grammar gives \textbf{\pa{ਹੋ}} as the stem of \textbf{\pa{ਹੋਣਾ}} and lists \textbf{\pa{ਹੁੰਦਾ}} as Present II. Punjabi University lesson 13 describes the agreeing \emph{-dā} construction with a present auxiliary, and its verb reference identifies \textbf{\pa{ਹੈ}} as the third-person singular present auxiliary. Sources: \href{https://www.learnpunjabi.org/statistics.html}{Punjabi University lexical statistics}, \href{https://pt.learnpunjabi.org/assets/A\%20Reference\%20Grammar\_Final.pdf}{Punjabi University reference grammar}, \href{https://pt.learnpunjabi.org/lesson13.aspx?l=13}{Punjabi University lesson 13}, and \href{https://www.learnpunjabi.org/verb.html}{Punjabi University verb reference}.

Those sources support the modern form. Reuse the history already attached to \textbf{\pa{ਹੋਣਾ}}; do not invent a second etymology for an inflected form. This chapter teaches the habitual reading only, not generic “is,” a progressive, counterfactual “would,” a past form, or a feminine or plural form.
\end{cousinweb}

\subsection*{Your turn}
\begin{itemize}
\raggedright
  \item \emph{Hear:} \emph{hundaa}; YOU SAY: \textbf{\pa{ਹੁੰਦਾ}}
  \item \emph{Read it:} \textbf{\pa{ਹੁੰਦਾ}}
  \item \emph{Write it:} copy \textbf{\pa{ਹੁੰਦਾ}} once, checking the whole form
\end{itemize}

\begin{tcolorbox}[breakable,colback=teal!4,colframe=teal!35!black,title={Before you move on}]
Produce the masculine singular habitual form for “happens.”
\end{tcolorbox}
\section[\texorpdfstring{\pa{ਹੁੰਦਾ} (hundaa)}{ਹੁੰਦਾ (hundaa)}]{Retrieve the new habitual form on its own}

\label{lesson:PA-C207-hundaa-review}



\begin{quote}
\emph{Hear:} \emph{hundaa}; YOU SAY: \textbf{\pa{ਹੁੰਦਾ}}
\end{quote}

\subsection*{Your turn}
\begin{itemize}
\raggedright
  \item \emph{Hear:} \emph{happens, masculine singular habitual}; YOU SAY: \textbf{\pa{ਹੁੰਦਾ}}
  \item \emph{Read it:} \textbf{\pa{ਹੁੰਦਾ}}
  \item \emph{Write it:} cover and write \textbf{\pa{ਹੁੰਦਾ}}
\end{itemize}

\begin{grammarlens}[title={Grammar Lens}]
Keep one whole form and one bounded habitual meaning here. The familiar frame arrives only after this separate retrieval.
\end{grammarlens}

\begin{tcolorbox}[breakable,colback=teal!4,colframe=teal!35!black,title={Before you move on}]
Say and then write the whole form once from memory.
\end{tcolorbox}
\section[\texorpdfstring{\pa{ਇਹ} \pa{ਰੋਜ਼} \pa{ਹੁੰਦਾ} \pa{ਹੈ} (eh roz hundaa hai)}{ਇਹ ਰੋਜ਼ ਹੁੰਦਾ ਹੈ (eh roz hundaa hai)}]{Put the new form in one daily frame}

\label{lesson:PA-C207-daily-frame}



\begin{quote}
Retrieve the familiar words for “this,” “every day,” and the final auxiliary \textbf{\pa{ਹੈ}}, then retrieve the new habitual form.
\end{quote}

\begin{grammarlens}[title={Grammar Lens}]
\begin{quote}
\textbf{\pa{ਇਹ} \pa{ਰੋਜ਼} \pa{ਹੁੰਦਾ} \pa{ਹੈ}} — \emph{eh roz hundaa hai} — this happens every day
\end{quote}

For this one frame, put familiar \textbf{\pa{ਇਹ} \pa{ਰੋਜ਼}} before the new whole form \textbf{\pa{ਹੁੰਦਾ}}, then keep familiar \textbf{\pa{ਹੈ}} last. Nothing else in the habitual or agreement system is added.
\end{grammarlens}

\subsection*{Your turn}
\begin{itemize}
\raggedright
  \item \emph{Hear:} \emph{eh roz hundaa hai}; YOU SAY: \textbf{\pa{ਇਹ} \pa{ਰੋਜ਼} \pa{ਹੁੰਦਾ} \pa{ਹੈ}}
  \item \emph{Read it:} \textbf{\pa{ਇਹ} \pa{ਰੋਜ਼} \pa{ਹੁੰਦਾ} \pa{ਹੈ}}
  \item \emph{Write it:} cover and write \textbf{\pa{ਇਹ} \pa{ਰੋਜ਼} \pa{ਹੁੰਦਾ} \pa{ਹੈ}}
\end{itemize}

\begin{tcolorbox}[breakable,colback=teal!4,colframe=teal!35!black,title={Before you move on}]
Produce the complete “this happens every day” frame.
\end{tcolorbox}
\section[\texorpdfstring{\pa{ਇਹ} \pa{ਰੋਜ਼} \pa{ਹੁੰਦਾ} \pa{ਹੈ} (eh roz hundaa hai)}{ਇਹ ਰੋਜ਼ ਹੁੰਦਾ ਹੈ (eh roz hundaa hai)}]{A bounded this-happens-every-day payoff}

\label{lesson:PA-C207-payoff}



\begin{quote}
Retrieve the new habitual form alone before rebuilding its familiar frame.
\end{quote}

\subsection*{Your turn: four skills}
\begin{itemize}
\raggedright
  \item \emph{Hear:} \emph{eh roz hundaa hai}; YOU SAY: \textbf{\pa{ਇਹ} \pa{ਰੋਜ਼} \pa{ਹੁੰਦਾ} \pa{ਹੈ}}
  \item \emph{Read it:} \textbf{\pa{ਇਹ} \pa{ਰੋਜ਼} \pa{ਹੁੰਦਾ} \pa{ਹੈ}}
  \item \emph{Write it:} cover the line and write \textbf{\pa{ਇਹ} \pa{ਰੋਜ਼} \pa{ਹੁੰਦਾ} \pa{ਹੈ}}
\end{itemize}

\begin{tcolorbox}[breakable,colback=teal!4,colframe=teal!35!black,title={Before you move on}]
Say the new form, then produce the complete daily frame.

This is the thirtieth prerequisite tranche exposed by Punjabi A2 mock drafting. It does not unblock Punjabi A2 mock 1 by itself. It is not an A2 mock, not A2 readiness or attainment, and not external or human validation.
\end{tcolorbox}
